% SGA 6, Exposición X: fragmento español alineado con la fuente.
% Alcance: sección 5 completa.
% Fuente francesa: sga6_fr_r1/src/fr.tex, líneas 13134--13295.
% Las lecturas fuente dudosas se conservan y se consignan en el informe de control.

\section*{5. El isomorfismo \(\Pic(X)\cong\Gr^1(X)\)}

\subsection*{5.1}
Sea \(X\) un topos anillado cualquiera. Entonces se pueden definir el anillo \(K^\circ(X)\) y la filtración \((\Filt^i(X))_{i\in\mathbb Z}\) exactamente como para los esquemas. Del mismo modo, se define el grupo \(\Pic(X)\) como el grupo de las clases de isomorfismo de los \(\mathcal O_X\)-módulos invertibles. Se denota por
\begin{equation}\tag{5.1.1}
 i:\Pic(X)\longrightarrow K^\circ(X)
\end{equation}
la aplicación que asocia a la clase \(\cl_P(L)\in\Pic(X)\) de un \(\mathcal O_X\)-módulo invertible \(L\) el elemento \(\cl^\circ(L)\) de \(K^\circ(X)\). Esta aplicación es evidentemente un homomorfismo de \(\Pic(X)\) al grupo multiplicativo \(K^\circ(X)^*\) de los elementos invertibles de \(K^\circ(X)\). En adelante escribiremos multiplicativamente la ley de grupo de \(\Pic(X)\).

La aplicación \(i\) tiene un inverso por la izquierda definido por el funtor «~det~»: cuando \(E\) es un \(\mathcal O_X\)-módulo localmente libre de rango \(n\), se tiene
\[
 \det(E)=\bigwedge^n E.
\]
Si
\begin{equation}\tag{5.1.2}
0\longrightarrow E'\longrightarrow E\longrightarrow E''\longrightarrow0
\end{equation}
es una sucesión exacta de \(\mathcal O_X\)-módulos localmente libres, hay un isomorfismo canónico
\[
 \det(E)\xrightarrow{\sim}(\det(E'))\otimes_{\mathcal O_X}(\det(E''))
\]
que proviene de (6.2.9) y del hecho de que (5.1.2) es localmente escindida. Se obtiene, por tanto, un homomorfismo de grupos abelianos
\begin{equation}\tag{5.1.3}
 \det:K^\circ(X)\longrightarrow\Pic(X)
\end{equation}
que es el inverso por la izquierda buscado de \(i\). En particular, \(i\) es inyectiva, y en adelante identificaremos \(\Pic(X)\) con su imagen en \(K^\circ(X)^*\).

\subsection*{5.2}
Se define otra aplicación
\begin{equation}\tag{5.2.1}
 \delta:\Pic(X)\longrightarrow\Filt^1(X)
\end{equation}
poniendo, para todo \(L\in\Pic(X)\),
\[
 \delta(L)=1-L^{-1}.
\]

\textbf{Observaciones 5.2.2.}:
\begin{enumerate}[label=\alph*)]
\item Supongamos que \(X\) es el topos localmente anillado de los haces de Zariski sobre un esquema \(S\). Cuando \(D\) es un divisor de Cartier \(\geq0\) sobre \(\sisi{X}{S}\) y % Canon SGA-SRC-001337
\[
 L=\cl^\circ(\mathcal O_S(D)),
\]
se tiene
\begin{equation}\tag{5.2.3}
 \theta_S(\delta(L))=\cl_\bullet(\mathcal O_D)
\end{equation}
donde \(\theta_S:K^\circ(X)\to K_\bullet(S)\) es el homomorfismo canónico de 1.2.3. Esto resulta inmediatamente de la sucesión exacta
\[
0\longrightarrow\mathcal O_S(-D)\longrightarrow\mathcal O_S\longrightarrow\mathcal O_D\longrightarrow0.
\]

\item Bajo las hipótesis de 5.1 y 5.2 se tiene, para todo \(L\in\Pic(X)\),
\begin{equation}\tag{5.2.4}
 \delta(L)\equiv L-1\quad\operatorname{mod.}\Filt^2(X),
\end{equation}
lo que justifica la notación \(c_1(L)\) introducida más abajo. En efecto, se tiene
\[
 (L-1)-(1-L^{-1})=(L-1)(1-L^{-1})\in\Filt^2(X).
\]
\end{enumerate}

\subsection*{5.3}
La composición de \(\delta\) con el epimorfismo canónico
\[
 \Filt^1(X)\longrightarrow\Gr^1(X)
\]
da la «~primera clase de Chern~» sobre \(\Pic(X)\):
\begin{equation}\tag{5.3.1}
 c_1:\Pic(X)\longrightarrow\Gr^1(X).
\end{equation}

\textbf{Teorema 5.3.2.} La aplicación \(c_1\) (5.3.1) es un homomorfismo inyectivo, y el homomorfismo
\[
 \det:\Filt^1(X)\longrightarrow\Pic(X)
\]
induce, al pasar al cociente, un homomorfismo inverso por la izquierda de \(c_1\):
\[
 \alpha:\Gr^1(X)\longrightarrow\Pic(X).
\]
Si toda sección \(n\) del haz constante \(\mathbb Z_X\) está acotada (por ejemplo, si \(X\) es «~cuasicompacto~», o si el objeto final \(e\) de \(X\) es suma finita de objetos conexos), entonces \(c_1\) es un isomorfismo y \(\alpha\) es su inverso.

\emph{Demostración:} \(c_1\) es un homomorfismo, pues para todo \(L,L'\in\Pic(X)\) se tiene la ecuación
\[
 (1-(LL')^{-1})-(1-L^{-1})-(1-{L'}^{-1})
 =
 -(1-L^{-1})(1-{L'}^{-1})\in\Filt^2(X).
\]
Admitamos por el momento que
\[
 \det(\Filt^2(X))=\{1\}\in\Pic(X),
\]
es decir, que \(\alpha\) está bien definido. Se tiene entonces
\[
 \alpha\circ c_1=\id_{\Pic(X)},
\]
pues, por 5.2.2 b), se tiene:
\[
 \det(1-L^{-1})=(\det(L-1))=L
\]
para todo \(L\in\Pic(X)\). A continuación se demuestra que
\[
 c_1\circ\alpha=\id_{\Gr^1(X)}
\]
cuando toda sección \(n\) de \(\mathbb Z_X\) está acotada. Como \(\Filt^1(X)\) es el grupo abeliano engendrado por los elementos \(\cl^\circ(E)-n\), donde \(E\) es un \(\mathcal O_X\)-módulo localmente libre de rango \(n\), basta demostrar, teniendo en cuenta 5.2.2 b), que
\begin{equation}\tag{5.3.3}
 \det(\cl^\circ(E))-\cl^\circ(E)+\sisi{N}{n}-1\in\Filt^2(X). % Canon SGA-SRC-001338
\end{equation}
Como \(n\) está acotada, el objeto final \(e\) de \(X\) se descompone en una suma directa finita disjunta \(\coprod e_i\) de subobjetos, de tal manera que \(n\) es constante sobre cada subtopos \(X/e_i\) \((i=1,\ldots,p)\). Como el funtor \(Y\rightsquigarrow K^\circ(X/Y)\) \((Y\in\operatorname{Ob}(X))\) transforma sumas directas finitas en productos directos finitos, se puede suponer que \(\underline n\) es constante sobre \(X\), es decir, \(\underline n=n\in\mathbb Z\). Poniendo \(E=\cl^\circ(\underline E)\), se obtiene (cf. Exp. V):
\[
 \det(E)=\lambda^n(E)=\gamma^n((E-n)+1)
 =
 \sum_{0\leq i\leq n}\gamma^i(E-n),
\]
lo que demuestra (5.3.3), pues el primer miembro es igual a
\[
 \sum_{2\leq i\leq n}\gamma^i(E-n).
\]

Queda por demostrar el

\textbf{Lema 5.3.4.} Para todo topos anillado \(X\) y para todo \(x\in\Filt^2(X)\), se tiene
\[
 \det(x)=1.
\]

\emph{Demostración:} Sea \(\epsilon:K^\circ(X)\to\Gamma(\mathbb Z)\) el homomorfismo de aumentación. Entonces (6.2.8) y (6.5.8), más abajo, dan por linealidad la fórmula
\begin{equation}\tag{5.3.5}
 \det(xy)=(\det(x))^{\epsilon(y)}(\det(y))^{\epsilon(x)}
 \qquad(x,y\in K^\circ(X)).
\end{equation}
Por tanto, basta demostrar que para todo \(k>1\) se tiene
\[
 \det(\gamma^k(E-n))=1,
\]
donde, como antes, \(E=\cl^\circ(\underline E)\), \(n=\epsilon(E)\). Ahora bien, se tiene
\[
 \gamma^k(x)=\lambda^k(x+k-1)
\]
para todo \(x\in K^\circ(X)\), de modo que
\[
 \det(\gamma^k(E-n))
 =
 \det(\lambda^k(E-n+k-1))
 =
 \det\left(\sum_{p+q=k}\binom{\sisi{n}{-n}+k-1}{p}\lambda^q(E)\right). % Canon SGA-SRC-001340
\]
Pero (6.2.10) y (6.5.3) dan
\[
 \det(\lambda^q(E))=(\det(E))^{\binom{n-1}{q-1}}
\]
para todo \(q>0\), de donde
\[
 \det(\gamma^k(E-n))=(\det(E))^r,
\]
con
\[
 r=\sum_{\substack{p+q=k\\ q\geq1}}
 \binom{-n+k-1}{p}\binom{n-1}{q-1}.
\]
Ahora bien, \(r\) es una sección de \(\mathbb Z\), igual en todas partes al coeficiente de \(t^{k-1}\) en \((1+t)^{k-2}\), y por tanto nula. Esto concluye la demostración.
